\documentclass[11pt,letterpaper]{article}
\usepackage[T1]{fontenc}
\usepackage[utf8]{inputenc}
\usepackage[margin=0.88in,top=0.84in,bottom=0.83in,headheight=13pt,headsep=20pt,footskip=28pt]{geometry}
\usepackage{newpxtext}
\usepackage{amsmath,amsthm}
\usepackage{newpxmath}
\usepackage{microtype}
\usepackage{xcolor}
\definecolor{Forest}{HTML}{30392C}
\definecolor{Olive}{HTML}{687144}
\definecolor{Muted}{HTML}{65685F}
\definecolor{Sage}{HTML}{CBD0BC}
\definecolor{Pale}{HTML}{F2F4EB}
\usepackage{mathtools}
\usepackage{booktabs,tabularx,longtable,array}
\usepackage{enumitem}
\usepackage{titlesec}
\usepackage{fancyhdr}
\usepackage[most]{tcolorbox}
\usepackage{needspace}
\PassOptionsToPackage{obeyspaces}{url}
\usepackage{xurl}
\usepackage[colorlinks=true,linkcolor=Olive,citecolor=Olive,urlcolor=Olive,bookmarksnumbered=true,pdfencoding=auto,psdextra]{hyperref}
\hypersetup{pdftitle={Normal measures hidden in critical-sequence quotients},pdfauthor={Prepared for Vladimir Reshetnikov},pdfsubject={Normal traces, Prikry reconstruction, and consistency calibrations},pdfkeywords={ultraexacting, HOD, normal measure, Prikry forcing, Mitchell order, equiconsistency}}
\urlstyle{same}
\setlength{\parindent}{0pt}
\setlength{\parskip}{5.1pt plus 1pt minus .4pt}
\linespread{1.035}
\setlength{\emergencystretch}{2em}
\setcounter{tocdepth}{1}
\setcounter{secnumdepth}{2}
\titleformat{\section}{\Large\bfseries\color{Forest}}{\thesection}{.6em}{}
\titleformat{\subsection}{\large\bfseries\color{Forest}}{\thesubsection}{.55em}{}
\titleformat{\subsubsection}{\normalsize\bfseries\color{Forest}}{}{0pt}{}
\titlespacing*{\section}{0pt}{18pt plus 3pt minus 2pt}{8pt}
\titlespacing*{\subsection}{0pt}{12pt plus 2pt minus 1pt}{5pt}
\titlespacing*{\subsubsection}{0pt}{9pt}{3pt}
\setlist[itemize]{leftmargin=1.4em,itemsep=2pt,topsep=3pt,parsep=0pt}
\setlist[enumerate]{leftmargin=1.7em,itemsep=3pt,topsep=4pt,parsep=0pt}
\pagestyle{fancy}
\fancyhf{}
\fancyhead[L]{\sffamily\fontsize{9}{11}\selectfont\color{Muted}LARGE CARDINALS AT THE INCONSISTENCY FRONTIER}
\fancyhead[R]{\sffamily\fontsize{9}{11}\selectfont\color{Muted}CONTINUATION IV}
\fancyfoot[L]{\small\color{Muted}Normal traces and Prikry reconstruction}
\fancyfoot[R]{\small\color{Muted}\thepage}
\renewcommand{\headrulewidth}{.35pt}
\renewcommand{\footrulewidth}{0pt}
\fancypagestyle{plain}{\fancyhf{}\fancyfoot[L]{\small\color{Muted}Normal traces and Prikry reconstruction}\fancyfoot[R]{\small\color{Muted}\thepage}\renewcommand{\headrulewidth}{0pt}}
\tcbset{colback=Pale,colframe=Sage,colbacktitle=Sage,coltitle=Forest,fonttitle=\bfseries,boxrule=.5pt,arc=3pt,left=9pt,right=9pt,top=6pt,bottom=6pt,toptitle=3pt,bottomtitle=3pt,before skip=8pt,after skip=8pt,breakable}
\newtcolorbox{assessment}[1]{title={#1}}
\theoremstyle{definition}
\newtheorem{definition}{Definition}[section]
\newtheorem{theorem}[definition]{Theorem}
\newtheorem{proposition}[definition]{Proposition}
\newtheorem{lemma}[definition]{Lemma}
\newtheorem{corollary}[definition]{Corollary}
\newtheorem{remark}[definition]{Remark}
\newtheorem{question}[definition]{Question}
\newtheorem{inputthm}{Input}
\renewcommand{\theinputthm}{\Alph{inputthm}}
\numberwithin{equation}{section}
\newcommand{\ZFC}{\mathsf{ZFC}}
\newcommand{\ZF}{\mathsf{ZF}}
\newcommand{\AC}{\mathsf{AC}}
\newcommand{\GA}{\mathsf{GA}}
\newcommand{\HOD}{\mathrm{HOD}}
\newcommand{\OD}{\mathrm{OD}}
\newcommand{\CD}{\mathrm{CD}}
\newcommand{\HCD}{\mathrm{HCD}}
\newcommand{\UF}{\mathrm{UF}}
\newcommand{\Ex}{\operatorname{Ex}}
\newcommand{\UEx}{\operatorname{UEx}}
\newcommand{\Ext}{\operatorname{Ext}}
\newcommand{\SC}{\operatorname{SC}}
\newcommand{\CEx}{\operatorname{CEx}}
\newcommand{\REx}{\operatorname{REx}}
\newcommand{\Con}{\operatorname{Con}}
\newcommand{\crit}{\operatorname{crit}}
\newcommand{\cf}{\operatorname{cf}}
\newcommand{\ot}{\operatorname{ot}}
\newcommand{\ran}{\operatorname{ran}}
\newcommand{\rank}{\operatorname{rank}}
\newcommand{\lcm}{\operatorname{lcm}}
\newcommand{\Ord}{\mathrm{Ord}}
\newcommand{\Pow}{\mathcal P}
\newcommand{\id}{\mathrm{id}}
\newcommand{\restr}{\mathbin{\upharpoonright}}
\newcommand{\Izero}{\mathrm{I0}}
\newcommand{\Itwo}{\mathrm{I2}}
\newcommand{\Ithree}{\mathrm{I3}}
\newcommand{\Iwf}{\mathrm{I3}_{\mathrm{wf}}(0)}
\newcommand{\Sel}{\operatorname{Sel}}
\newcommand{\FF}{\mathcal F}
\newcommand{\PP}{\mathbb P}
\newcommand{\Clam}{\mathcal C_\lambda}
\newcommand{\Rig}{\operatorname{Rig}}
\newcommand{\Spec}{\operatorname{Spec}}
\newcommand{\ch}{\operatorname{ind}}
\newcommand{\CF}{\mathsf{CF}}
\newcommand{\src}[1]{\unskip\ {\normalfont\small\sffamily\color{Olive}[#1]}\ }
\newcommand{\R}[1]{\textsf{R#1}}
\newcolumntype{Y}{>{\raggedright\arraybackslash}X}
\newcolumntype{P}[1]{>{\raggedright\arraybackslash}p{#1}}
\newcolumntype{C}{>{\centering\arraybackslash}p{.034\textwidth}}
\newcommand{\yes}{$\bullet$}
\newcommand{\prt}{$\circ$}
\newcommand{\arxiv}[1]{\href{https://arxiv.org/abs/#1}{arXiv:#1}}
\newcommand{\doi}[1]{\href{https://doi.org/#1}{doi:\nolinkurl{#1}}}


\newcommand{\tc}{\operatorname{tc}}
\newcommand{\dom}{\operatorname{dom}}
\newcommand{\Ult}{\operatorname{Ult}}
\newcommand{\Meas}{\operatorname{Meas}}
\newcommand{\NT}{\mathsf{NT}}
\newcommand{\NTP}{\mathsf{NT}^{+}}
\newcommand{\normal}{\mathsf{NM}}
\newcommand{\ODp}{\OD_{\{p\}}}
\newcommand{\HODp}{\HOD_{\{p\}}}
\newcommand{\tail}[2]{#1\mathbin{/}#2}
\newcommand{\fin}{\mathrm{fin}}
\newcommand{\seq}[1]{\langle #1\rangle}
\newcommand{\force}{\Vdash}
\newcommand{\concat}{\mathbin{{}^{\frown}}}
\begin{document}
\thispagestyle{empty}
{\sffamily\bfseries\small\color{Olive}SET THEORY\quad /\quad RESEARCH CONTINUATION IV}

\vspace{13pt}
{\fontsize{28}{31}\selectfont\bfseries\color{Forest}Normal measures hidden in\\ critical-sequence quotients\par}
\vspace{10pt}
{\fontsize{14}{18}\selectfont A Prikry factorization of ultraexacting witnesses\\ and two exact consistency calibrations\par}
\vspace{14pt}
{\color{Muted}Prepared for Vladimir Reshetnikov\hfill 18 September 2026\par}
\vspace{3pt}
{\color{Sage}\hrule height .6pt}
\vspace{12pt}

\textbf{Abstract.}
The supplied synthesis obtains regularity in parameterized HOD models from
rigid classes of cofinal $\omega$-sets modulo finite difference. We develop a
stronger conclusion for a critical-sequence class. At an ultraexacting
cardinal $\lambda$, one can choose a bijection $b:\lambda\to V_\lambda$ and
$q=[a]$ so that the canonical eventual-containment trace
\[
 U_{b,q}=\{A\in\Pow(\lambda)\cap\HOD_{\{b,q\}}:
                   a\subseteq^* A\}
\]
is a normal measure \emph{inside} $N=\HOD_{\{b,q\}}$, concentrating on the
cardinals below $\lambda$ that are measurable in $N$. Every representative
of $q$ is Prikry-generic over $N$ for this measure. Thus the ambient
singularity of $\lambda$ admits a canonical Prikry subextension with ground
model agreeing with the ambient universe through $V_\lambda$.

The isolated existence of such a normal trace is equiconsistent with a
measurable cardinal. Requiring the trace to concentrate on measurables is
equiconsistent with a cardinal of Mitchell order at least two. The upper
bounds use a quotient-relative cone-homogeneity argument for ordinary
Prikry forcing; the lower bounds use the explicitly defined inner model.
These calibrations concern the trace principles, not ultraexactingness
itself. Detailed conventional proofs are given, including the required
Prikry genericity and homogeneity arguments.

\begin{assessment}{The principal advance}
For a suitable quotient-parameter HOD model, the change is
\[
 \text{regularity}\ \longrightarrow\ \text{normal measure}
 \ \longrightarrow\ \text{Prikry genericity}.
\]
A further, independent conclusion identifies the precise consistency
strength of the two isolated trace principles. No inconsistency of
$\ZFC+\Izero$ or of ultraexacting cardinals is asserted.
\end{assessment}

\textbf{Status.} Research manuscript with full English proofs; not
refereed or formalized in Lean. The new assertions are presented as
continuations of the supplied synthesis, not as a claim of established
literature priority. Published inputs and the limits of the conclusions
are separated explicitly from the arguments developed here.

\newpage
\tableofcontents
\newpage

\section{Statements and their relation to the supplied research}\label{sec:results}

We work in $\ZFC$ unless an inner model is explicitly named. A \emph{normal
measure on $\kappa$} means a nonprincipal, $\kappa$-complete ultrafilter on
$\Pow(\kappa)$ with the regressive-function normality property. The words
``on $\Pow(\kappa)$'' and the completeness quantifier are always evaluated
in the model under discussion.

For an uncountable cardinal $\lambda$ with $\cf^V(\lambda)=\omega$, set
\[
 D_\lambda=\{a\subseteq\lambda:\ot(a)=\omega\text{ and }\sup a=\lambda\},
 \qquad Q_\lambda=D_\lambda/{=^*},
\]
where $a=^*d$ means that $a\triangle d$ is finite. Write $[a]$ for the
corresponding \emph{set} of representatives. In particular, a parameter
$q\in Q_\lambda$ is one set parameter; it does not give its individual
members as freely available parameters.

For a bijection $b:\lambda\to V_\lambda$ and $q\in Q_\lambda$, define
\begin{align}
 N_{b,q}&=\HOD_{\{b,q\}}^V,\label{eq:N}\\
 U_{b,q}&=\{A\in\Pow(\lambda)\cap N_{b,q}:
                  (\exists a\in q)\ a\subseteq^* A\}.
                  \label{eq:trace}
\end{align}
Here $a\subseteq^* A$ means $a\setminus A$ is finite. The existential
quantifier in \eqref{eq:trace} can equally well be universal over $q$.
The fixed-parameter convention and the fact that $N_{b,q}$ satisfies
Choice are proved in Section~\ref{sec:hod}.

\begin{definition}[The two trace principles]\label{def:NT}
The assertion $\NT(\lambda,b,q)$ says that $\lambda$ is an uncountable
cardinal of ambient cofinality $\omega$, that $b:\lambda\to V_\lambda$ is
a bijection, that $q\in Q_\lambda$, and that
\[
 N_{b,q}\models\text{``$U_{b,q}$ is a normal measure on $\lambda$.''}
\]
The assertion $\NTP(\lambda,b,q)$ adds
\[
 \Meas_\lambda^{N_{b,q}}
 =\{\alpha<\lambda:N_{b,q}\models\text{``$\alpha$ is measurable''}\}
 \in U_{b,q}.
\]
Let $T_{\NT}$ and $T_{\NTP}$ denote $\ZFC$ together with the existence of
such a triple for the respective principle.
\end{definition}

These are first-order assertions after expanding the definable HOD
predicate and relativizing the finitely many formulas about measures.
The definition does \emph{not} assert the existence of an ambient normal
measure on the singular cardinal $\lambda$.

\begin{theorem}[Ultraexacting normal trace]\label{thm:main}
If $\lambda$ is ultraexacting, there are $b$ and $q$ such that
$\NTP(\lambda,b,q)$ holds. The class $q$ can be the finite-difference
class of the critical sequence of an ultraexacting witness at a
sufficiently correct height. For $N=N_{b,q}$ and $U=U_{b,q}$:
\begin{enumerate}
\item $V_\lambda^N=V_\lambda$ and $N$ sees $\lambda$ as measurable, with
a normal measure concentrating on measurables.
\item Every $a\in q$, in its increasing enumeration, is
$\PP_U^N$-generic over $N$.
\item $N[a]$ is an inner model of $V$ satisfying
\[
  V_\lambda^{N[a]}=V_\lambda,
  \qquad \cf^{N[a]}(\lambda)=\omega.
\]
Moreover $N[a]=N[d]$ for every $a,d\in q$.
\end{enumerate}
\end{theorem}

The measure construction is proved in Section~\ref{sec:normal}; the
Prikry conclusions are proved in Section~\ref{sec:prikry}. No assumption
about iterability of an infinite iteration of the witness is used.

\begin{theorem}[Exact consistency calibrations]\label{thm:calibrations}
The following equivalences of consistency statements hold:
\begin{align}
 \Con(T_{\NT})
  &\ \longleftrightarrow\
    \Con(\ZFC+\text{there is a measurable cardinal}),
       \label{eq:conNT}\\
 \Con(T_{\NTP})
  &\ \longleftrightarrow\
    \Con(\ZFC+\exists\kappa\,o(\kappa)\ge2).
       \label{eq:conNTplus}
\end{align}
In the second line, $o(\kappa)\ge2$ may equivalently be replaced by:
there is a normal measure on $\kappa$ concentrating on measurable
cardinals below $\kappa$.
\end{theorem}

The upper bounds are actual forcing constructions, not appeals to a
large-cardinal implication diagram. The proofs are in
Sections~\ref{sec:homogeneity}--\ref{sec:calibration}.

\subsection{What has, and has not, been extended}

The attached \emph{Synthesis} proves that certain $q$ make $\lambda$
regular in $\HOD_{V_\lambda\cup\{q\}}$ and develops sharp-width results
for definable sections of $Q_\lambda$ \cite{source-synthesis}. The
present construction uses a special critical-sequence class and a fixed
bijection $b$. Its model $\HOD_{\{b,q\}}$ contains the former
parameter-HOD model: each individual parameter in $V_\lambda$ can be
recovered from $b$ and an ordinal. This is a stronger conclusion for a
specified model, but is \emph{not} a claim that every rigid class carries
the measure in \eqref{eq:trace}.

The three assertions in \eqref{eq:conNT}, \eqref{eq:conNTplus}, and
``there is an ultraexacting cardinal'' must be kept distinct. The
published equiconsistency of ultraexactingness with $\Izero$ remains
unchanged \cite[Theorem~4.5]{abgl}. We isolate consequences of that much
stronger hypothesis and determine their own exact strengths.

The fixed-witness and canonical-choice methods have precedents in the
supplied reports and in the work on exacting cardinals. The continuation
here is the conversion of their eventual behavior into normality, the
Prikry reconstruction, and the converse quotient-homogeneity
construction. Classical Prikry facts are identified and re-proved rather
than counted as new results.

\section{Fixed-parameter HOD and the canonical trace}\label{sec:hod}

\subsection{A necessary parameter convention}

For a fixed finite tuple $p$ of sets, $\ODp$ denotes sets definable from
$p$ and finitely many ordinal parameters. The tuple may be encoded as
one set. Define
\[
 \HODp=\{x:\tc(\{x\})\subseteq\ODp\}.
\]
This differs from allowing \emph{arbitrary members} of a possibly
unwellorderable set as parameters. All HOD models with braces in this
paper use only the displayed finite tuple.

\begin{lemma}[Fixed-parameter inner models]\label{lem:hod}
For every fixed finite tuple $p$, $\HODp$ is a definable transitive inner
model of $\ZFC$. The class $\ODp$ has a uniformly definable, set-like
canonical well-order. A set of ordinals is in $\HODp$ if and only if it is
in $\ODp$.
\end{lemma}

\begin{proof}
A definition can be coded by a rank height, a formula number, and a finite
tuple of ordinal parameters. One uses truth in a set structure
$V_\theta$ containing $p$, rather than a truth predicate for $V$.
Reflection identifies the resulting definability notion with ordinal
definability from $p$. Order the codes first by a bound on their height
and ordinal entries, and then by a fixed well-order within that bound.
Assign each definable object its first code. This produces a uniform
set-like well-order of $\ODp$; in particular every nonempty set of
$\ODp$ objects has a canonical first member.

The hereditary definition gives transitivity, and all ordinals belong
to the class. The usual HOD closure argument applies with $p$ held
fixed. For example, the internal power set of $x\in\HODp$ is
$\Pow(x)\cap\HODp$. This is definable from $p$ and a definition of $x$;
its members and all their descendants are in $\ODp$. It therefore
belongs to $\HODp$. Separation and Replacement follow by relativizing
a formula to the definable class $\HODp$ and substituting definitions
for its finitely many parameters. Pairing, Union, Infinity, Foundation,
and Extensionality are handled in the same way or by transitivity.

Restrict the canonical order to any set in $\HODp$. The resulting
well-order is itself hereditarily definable from $p$, since its ordered
pairs are made from hereditarily definable elements. Thus it belongs to
the model and proves Choice there. Finally, the descendants of a set
of ordinals, other than the set itself, are ordinals. They are already
ordinal definable, proving the last assertion.
\end{proof}

The lemma does not say $p\in\HODp$. For example, $q$ will be definable
from itself but will have no representative in the inner model. Its
failure of hereditary definability is essential.

\begin{lemma}[A rank-bijection supplies the lower universe]\label{lem:b-rank}
Suppose $b:\lambda\to V_\lambda$ is a bijection and $q$ is any set. Then
$N=\HOD_{\{b,q\}}$ contains every member of $V_\lambda$,
$V_\lambda^N=V_\lambda$, and $b\in N$.
\end{lemma}

\begin{proof}
For $x\in V_\lambda$, choose the ordinal $\xi<\lambda$ with $b(\xi)=x$.
Then $x$ is definable from $b$ and $\xi$. Every member of its transitive
closure also lies in $V_\lambda$ and has the same kind of definition.
Consequently $x\in N$. Ranks of sets in transitive inner models are
absolute, so this proves the equality of rank segments. The set $b$
itself is definable from the displayed parameters. Its graph elements
and all their descendants are in $V_\lambda$ when $\lambda$ is a limit
cardinal, as in all applications here. Hence $b\in N$. More generally,
the same conclusion follows directly by hereditary definability of
ordered pairs of ordinals and values of $b$.
\end{proof}

\begin{lemma}[The trace is an internal set]\label{lem:trace-set}
For a triple as in Definition~\ref{def:NT}, without assuming normality,
$U_{b,q}\in N_{b,q}$. It is a proper filter on $\Pow(\lambda)^{N_{b,q}}$
containing every final segment of $\lambda$.
\end{lemma}

\begin{proof}
Finite modification does not change eventual containment, so the
condition in \eqref{eq:trace} is independent of the representative.
The set $U_{b,q}$ is definable from $b,q$ and $\lambda$, because
$N_{b,q}$ is a definable class. Every member $A$ is in $N_{b,q}$ and
hence is hereditarily definable from $b,q$. This shows that the trace
itself is hereditarily definable and belongs to $N_{b,q}$.

The empty set does not eventually contain a cofinal $\omega$-set. The
whole cardinal does. Eventual containment is upward closed and is
preserved by intersections of two sets, giving the filter properties.
An increasing $\omega$-sequence cofinal in $\lambda$ has only finitely
many terms below any fixed $\beta<\lambda$: otherwise all its terms
would occur below $\beta$, contradicting cofinality. Thus every final
segment belongs to the trace.
\end{proof}

Normality is a substantial additional assertion. An arbitrary cofinal
sequence can alternate between a definable set and its complement, so
Lemma~\ref{lem:trace-set} alone does not even give an ultrafilter.

\Needspace{9\baselineskip}
\section{Correct heights, critical sequences, and fixed parameters}\label{sec:fixed}

\subsection{The exact amount of reflection needed}

The notions originate in \cite{abl}. We use the following form of the
definition from \cite[Section~2]{abgl}. At an ultraexacting $\lambda$, for arbitrarily
high prescribed sufficiently correct cardinal heights $\zeta$, there
are
\begin{equation}\label{eq:witness}
 X\prec V_\zeta,\qquad V_\lambda\cup\{\lambda\}\subseteq X,
 \qquad j:X\longrightarrow V_\zeta
\end{equation}
with $j(\lambda)=\lambda$, $\crit(j)<\lambda$, and
$e=j\restr V_\lambda\in X$. Here $e:V_\lambda\to V_\lambda$ is
elementary. We choose the height \emph{before} choosing this witness.

\begin{lemma}[Uniform canonical-counterexample principle]\label{lem:least}
Fix finitely many first-order counterexample predicates, with a finite
parameter tuple $p$ and a cardinal parameter $\lambda$. Assume their
candidate objects have rank below $\lambda+\omega$. There is a finite
reflection requirement on $V_\zeta$, uniform in $p$, with the following
consequence. If \eqref{eq:witness} holds, $p\in X$, and $j(p)=p$, then
any nonempty class of $\ODp$ candidates for one of these predicates has
a candidate $x\in X$ with $j(x)=x$.
\end{lemma}

\begin{proof}
For each predicate, write a single formula saying that $x$ is its
canonical $\ODp$-least counterexample. Definability and the canonical
order in Lemma~\ref{lem:hod} are first-order definable uniformly in
$p$. The finitely many resulting formulas, their subformulas needed for
reflection, and existence and uniqueness of the selected object form a
finite list. Choose a cardinal $\zeta>\lambda+\omega$ at which they are
correct for all parameters in $V_\zeta$. Equivalently, membership in
$C^{(m)}$ for one sufficiently large fixed finite $m$ suffices. Such
heights are supplied by L\'evy reflection.

If counterexamples exist in $V$, their canonical first member $x$ has
the stipulated bounded rank, so $x\in V_\zeta$. Correctness says that
$V_\zeta$ selects exactly this same $x$. Because $p,\lambda\in X$ and
$X\prec V_\zeta$, the unique selected object belongs to $X$.
Elementarity of $j$, together with $j(p)=p$ and $j(\lambda)=\lambda$,
says that $j(x)$ is the same unique object of $V_\zeta$. Thus $j(x)=x$.
\end{proof}

There is no assumption that the ordinal parameters in an arbitrary
$\ODp$ definition are fixed by $j$. Such an assumption would be false.
The canonical least object is fixed because it is uniquely definable
from the \emph{fixed tuple}, not because its original defining code is
pointwise fixed. Uniform reflection also avoids choosing a height only
after the witness has supplied $b$ and $q$.

\subsection{The critical sequence and its lower-rank measures}

\begin{lemma}[Critical sequence facts]\label{lem:critical}
For a witness \eqref{eq:witness} let
\[
 \kappa=\crit(j)=\crit(e),\qquad
 \kappa_n=e^n(\kappa),\qquad a=\{\kappa_n:n<\omega\}.
\]
Then every $\kappa_n$ is measurable in $V$,
$\sup_{n<\omega}\kappa_n=\lambda$, and $\lambda$ is a strong limit
cardinal of cofinality $\omega$. The enumeration of $a$ and the set $a$
are in $X$, and
\[
 j(a)=a\setminus\{\kappa\}.
\]
Consequently $q=[a]\in X$ and $j(q)=q$.
\end{lemma}

\begin{proof}
First recall the derived-measure argument. Set
\[
 U_0=\{A\subseteq\kappa:\kappa\in e(A)\}.
\]
All relevant subsets, functions, and sequences below $\kappa$ have
rank below $\lambda$, so $e$ acts on them. It follows immediately from
complements and finite intersections that $U_0$ is an ultrafilter.
Every singleton is null, since ordinals below $\kappa$ are fixed. If
$\mu<\kappa$ and $\seq{A_i:i<\mu}$ consists of members of $U_0$, then
$e$ fixes the index set and sends its intersection to the intersection
of the images. Thus $\kappa$ belongs to that image and the intersection
lies in $U_0$.

For a regressive $f:S\to\kappa$ with $S\in U_0$, we have
$e(f)(\kappa)=\beta<\kappa$. The ordinal $\beta$ is fixed, so
$\kappa\in e(\{\alpha\in S:f(\alpha)=\beta\})$. This proves normality.
The critical point is uncountable, since finite ordinals and $\omega$
are fixed. Hence $U_0$ witnesses measurability. The statement that it is
a normal measure is absolute between $V_\lambda$ and $V$, all its
quantifiers ranging over sets of bounded rank here. Applying $e^n$ to
$U_0$ therefore gives a normal measure on $\kappa_n$ in $V$.

We use the classical rank form of Kunen's inconsistency: there is no
nontrivial elementary self-embedding of $V_{\delta+2}$ in $\ZFC$
\cite{kunen}. If $\delta=\sup_n\kappa_n<\lambda$, the critical sequence
has bounded rank in $V_\lambda$, and $e(\delta)=\delta$. Restricting
$e$ to $V_{\delta+2}$ would give just such a forbidden embedding. Thus
$\delta=\lambda$. Measurable cardinals are strongly inaccessible, so
$\lambda$ is a countable increasing limit of strong limits and is itself
strong limit. For completeness, regularity of a measurable follows by
ruling out a partition into fewer than $\kappa$ bounded null sets;
strong limit follows by taking a purported injection
$\kappa\to\Pow(\theta)$ for $\theta<\kappa$ and intersecting the
measure-one decisions of all $\theta$ membership coordinates. The
resulting measure-one fiber would have at most one element.

The recursion defining $\seq{\kappa_n:n<\omega}$ is a set recursion in
$V_\zeta$ from $e$. Since $e\in X$, elementarity places its unique
value in $X$. A function with domain $\omega$ is mapped pointwise by
$j$, and $j(\kappa_n)=e(\kappa_n)=\kappa_{n+1}$. This proves the stated
formula for $j(a)$. Forming the finite-difference class is definable
from $a,\lambda$, and is correctly computed at this height. Deleting
one element does not change the class, so $q\in X$ is fixed.
\end{proof}

Only the explicitly stated classical Kunen theorem is imported in this
argument. No assertion about an iterable direct limit of the iterates
of $e$ is required.

\begin{lemma}[A fixed bijection]\label{lem:fixed-b}
The witness in Lemma~\ref{lem:critical} supplies a bijection
$b:\lambda\to V_\lambda$ with $b\in X$ and $j(b)=b$.
\end{lemma}

\begin{proof}
As $\kappa$ is inaccessible, choose a bijection
$b_0:\kappa\to V_\kappa$. Its rank is below $\lambda$, so
$b_0\in V_\lambda\subseteq X$. Define
\[
 b_n=e^n(b_0):\kappa_n\longrightarrow V_{\kappa_n}.
\]
Every ordinal below $\kappa$ and every member of $V_\kappa$ is fixed by
$e$. Therefore $b_0\subseteq e(b_0)$. Applying $e$ repeatedly gives
$b_n\subseteq b_{n+1}$. Consequently
\[
 b=\bigcup_{n<\omega} b_n
\]
is a bijection from $\lambda$ to $V_\lambda$.

The sequence of the $b_n$ is uniquely obtained by recursion from $e$
and $b_0$, both in $X$, so the sequence and its union belong to $X$.
Each $b_n\in V_\lambda$, hence $j(b_n)=e(b_n)=b_{n+1}$. Mapping the
countable sequence pointwise and taking unions gives
\[
 j(b)=\bigcup_{n<\omega} b_{n+1}=b.
\]
This is the bijection version of the fixed-well-order device also used
in \cite[Proposition~3.9]{abgl} and in the supplied synthesis.
\end{proof}

\section{From eventual decisions to a normal measure}\label{sec:normal}

Fix the height in Lemma~\ref{lem:least} for the two predicates used
below, then choose the ultraexacting witness and its $a,q,b$. Write
$p=(b,q)$, $N=\HODp$, and $U=U_{b,q}$. Every parameter in $p$ is fixed
by $j$.

\begin{lemma}[Eventual decisions]\label{lem:decisions}
For every $A\subseteq\lambda$ in $\ODp$, either $a\subseteq^* A$ or
$a\subseteq^*(\lambda\setminus A)$.
\end{lemma}

\begin{proof}
Suppose otherwise. A counterexample is an $\ODp$ subset of $\lambda$
that splits a representative of $q$ into two infinite pieces. This
property is independent of the representative, is first-order in
$p,\lambda$, and has candidates of bounded rank. By
Lemma~\ref{lem:least} there is such a counterexample $A_0\in X$ with
$j(A_0)=A_0$.

For every $n<\omega$, elementarity gives
\[
 \kappa_n\in A_0\quad\Longleftrightarrow\quad
 j(\kappa_n)=\kappa_{n+1}\in j(A_0)=A_0.
\]
Membership in $A_0$ is therefore constant along the entire critical
sequence. Thus either $a\subseteq A_0$ or $a\cap A_0=\varnothing$,
contradicting the choice of $A_0$.
\end{proof}

\begin{lemma}[Eventual pressing down]\label{lem:press}
Let $f\in\ODp$ have domain $S\subseteq\lambda\setminus\{0\}$ and
satisfy $f(\alpha)<\alpha$ for $\alpha\in S$. If $a\subseteq^*S$, then
$f$ is eventually constant along $a$.
\end{lemma}

\begin{proof}
Failure is again a first-order counterexample predicate in $p,\lambda$:
the graph is a set of ordinal pairs of rank below $\lambda+\omega$;
its domain eventually contains a representative of $q$; and there is
no ordinal $\beta$ with that representative eventually contained in the
fiber $f^{-1}\{\beta\}$. Select a canonical counterexample $f_0\in X$
with $j(f_0)=f_0$ by Lemma~\ref{lem:least}.

The set $S_0=\dom(f_0)$ is fixed. Membership of $\kappa_n$ in $S_0$ is
constant in $n$ by the same calculation as in
Lemma~\ref{lem:decisions}. Since membership holds eventually, it holds
for every $n$, including $n=0$. Put $\beta=f_0(\kappa)$. Regressiveness
gives $\beta<\kappa$, so $j(\beta)=\beta$. Elementarity of function
evaluation now yields, successively,
\[
 f_0(\kappa_{n+1})=j\bigl(f_0(\kappa_n)\bigr)=\beta.
\]
The function is constant on all of $a$, a contradiction.
\end{proof}

\begin{proposition}[Internal normality and completeness]\label{prop:normal}
The set $U$ belongs to $N$ and $N$ regards it as a normal measure on
$\lambda$.
\end{proposition}

\begin{proof}
Membership $U\in N$ and the proper filter properties were proved in
Lemma~\ref{lem:trace-set}. Every $A\in\Pow(\lambda)^N$ is in $\ODp$,
so Lemma~\ref{lem:decisions} decides it or its complement. These two
alternatives cannot both occur, because $a$ is infinite. Thus $U$ is an
ultrafilter on the full internal power set. It contains all final
segments and no singleton, so it is nonprincipal.

If $S\in U$ and $f\in N$ is regressive on $S\setminus\{0\}$, then
$f\in\ODp$ and $a\subseteq^* S$. Lemma~\ref{lem:press} gives an
ordinal $\beta$ such that $a\subseteq^* f^{-1}\{\beta\}$. The fiber is
in $N$ and in $U$. This is the regressive-function normality property
for \emph{all} such $f$ in $N$.

We verify $\lambda$-completeness rather than presupposing it. Let
$\seq{A_i:i<\mu}\in N$, where $\mu<\lambda$ and each $A_i\in U$.
Its intersection belongs to $N$. Suppose the intersection were not in
$U$. The ultrafilter property would put its complement in $U$. After
intersecting with the final segment above $\mu$, obtain $B\in U$ and
define in $N$
\[
 f(\alpha)=\min\{i<\mu:\alpha\notin A_i\}\qquad(\alpha\in B).
\]
This is regressive. Normality gives an $i<\mu$ whose fiber is in $U$.
But that fiber is disjoint from $A_i\in U$, contradicting the proper
filter property. The case $\mu=0$ is immediate. This proves
completeness for every sequence of length less than $\lambda$ that
belongs to $N$.
\end{proof}

\begin{proposition}[Concentration on measurables]\label{prop:concentration}
For $M_\lambda=\Meas_\lambda^N$, we have $a\subseteq M_\lambda$ and
$M_\lambda\in U$. In particular $\NTP(\lambda,b,q)$ holds.
\end{proposition}

\begin{proof}
The model $N$ contains $V_\lambda$ by Lemma~\ref{lem:b-rank}. For each
$n$, the normal measure on $\kappa_n$ obtained in
Lemma~\ref{lem:critical} has rank below $\lambda$, and hence belongs to
$N$. All subsets of $\kappa_n$ and all functions needed to test its
normality and completeness also have rank below $\lambda$. Thus the
same measure witnesses measurability in $N$.

The set $M_\lambda$ is defined by Separation inside $N$ and consists of
ordinals, so it is an eligible argument of the trace. It contains every
term of $a$, proving $M_\lambda\in U$ directly from the definition.
Together with Proposition~\ref{prop:normal}, this establishes the
measure part of Theorem~\ref{thm:main}.
\end{proof}

\begin{assessment}{Where the additional strength enters}
Eventual decisions give an ultrafilter. Eventual pressing down gives
normality. A regressive ``least failed coordinate'' then supplies
internal $\lambda$-completeness. Merely knowing that $\lambda$ is
regular in a parameterized HOD model would not supply any of these
three steps.
\end{assessment}

\section{Prikry reconstruction, with the genericity proof}\label{sec:prikry}

The results in this section about normal measures and Prikry forcing are
classical. The genericity criterion is due to Mathias \cite{mathias};
see also the formulation in \cite{merimovich}. We include proofs because
the fact that our inner model need not be countable makes a mere
``meet its dense sets one by one'' argument inappropriate.

\subsection{The forcing and finite-stem diagonal intersections}

Temporarily work inside a transitive model $M\models\ZFC$ with a normal
measure $W$ on $\kappa$. A condition in $\PP_W$ is a pair $(s,A)$, where
$s$ is a finite increasing sequence of ordinals below $\kappa$,
$A\in W$, and every element of $A$ exceeds every entry of $s$. Empty
stems are allowed. The order is
\[
 (t,B)\le(s,A)
 \quad\Longleftrightarrow\quad
 \begin{cases}
  s\text{ is an initial segment of }t,\\
  \text{all new entries of }t\text{ lie in }A,\\
  B\subseteq A.
 \end{cases}
\]
A direct extension, written $\le^*$, has the same stem. For a finite
sequence $t$, write $A/t$ for the elements of $A$ above every entry of
$t$; thus $A/\varnothing=A$.

\begin{lemma}[Diagonalization over finite stems]\label{lem:diag}
If $\seq{A_s:s\in[\kappa]^{<\omega}}\in M$ and each $A_s\in W$, then
\[
 \Delta_{\fin} A_s
 =\{\beta<\kappa:(\forall s\in[\beta]^{<\omega})\ \beta\in A_s\}
 \in W.
\]
Finite subsets and their increasing enumerations are identified here.
\end{lemma}

\begin{proof}
All choices and functions in this proof are taken in $M$. Suppose the
complement were in $W$. Intersect it with $A_\varnothing$ and with
$\kappa\setminus\{0\}$. For each remaining $\beta$, choose a nonempty
finite $s_\beta\subseteq\beta$ with $\beta\notin A_{s_\beta}$. The
function $\beta\mapsto\max(s_\beta)$ is regressive. Normality makes it
constant, with value $\gamma$, on a set in $W$.

There are fewer than $\kappa$ finite subsets of $\gamma+1$. By
$\kappa$-completeness and the ultrafilter property, on a further set in
$W$ the whole finite set $s_\beta$ has a constant value $s$. That set
is disjoint from $A_s\in W$, a contradiction. The intersection with
$A_\varnothing$ also accounts for the empty-stem case.
\end{proof}

\begin{lemma}[Finite homogeneity]\label{lem:ramsey}
For every $n<\omega$, every $\rho<\kappa$, and every coloring
$h:[\kappa]^n\to\rho$ in $M$, there is $H\in W$ homogeneous for $h$.
\end{lemma}

\begin{proof}
For $n=0$ there is only one input. For $n=1$, a partition into fewer
than $\kappa$ color classes has exactly one class in $W$. Suppose the
result holds for $n$, and consider a coloring of $(n+1)$-element sets.
For each $s\in[\kappa]^n$, there is a color $i_s<\rho$ such that
\[
 B_s=\{\beta>\max s:h(s\cup\{\beta\})=i_s\}\in W.
\]
For $n=0$ the maximum restriction is omitted. Existence of $i_s$ again
uses $\kappa$-completeness. Apply the induction hypothesis to the
coloring $s\mapsto i_s$ and obtain a homogeneous $H_0\in W$.
Intersect $H_0$ with the finite-stem diagonal intersection of the $B_s$
(assign $\kappa$ to unused stem lengths). On the resulting set $H$,
every increasing $(n+1)$-tuple has its last entry in the appropriate
$B_s$ and its first $n$ entries in $H_0$. Its color is therefore the
single color fixed by the induction hypothesis.
\end{proof}

\begin{lemma}[Strong Prikry lemma]\label{lem:strong-prikry}
If $D\in M$ is dense open in $\PP_W$ and $(s,A_0)$ is a condition,
there are $A\subseteq A_0$ in $W$ and $n<\omega$ such that
\[
 (s\concat t,A/t)\in D\qquad\text{for every }t\in[A]^n.
\]
\end{lemma}

\begin{proof}
For every finite increasing $t$ from $A_0$, color $t$ good if there
exists $B\in W$ such that $(s\concat t,B)\in D$, and bad otherwise.
For each length separately, Lemma~\ref{lem:ramsey} makes this coloring
constant on a set in $W$. Intersect these countably many sets with
$A_0$. Since $\kappa$ is uncountable and $W$ is $\kappa$-complete,
the result $H$ still belongs to $W$.

By density, some $(s\concat t,B)\le(s,H)$ belongs to $D$. Thus the
homogeneous color at length $n=|t|$ is good. For each $u\in[H]^n$
choose a witness $B_u\in W$ with $(s\concat u,B_u)\in D$.
Apply Lemma~\ref{lem:diag} to these reservoirs, assigning $\kappa$ to
all unused stems, and let $A$ be the intersection of that diagonal set
with $H$. If $u\in[A]^n$ and $\beta\in A/u$, then
$u\subseteq\beta$, so $\beta\in B_u$. Hence $A/u\subseteq B_u$.
The condition $(s\concat u,A/u)$ extends the chosen condition in $D$,
and belongs to $D$ because $D$ is open. This includes $n=0$, using the
empty-stem reservoir in the diagonal intersection.
\end{proof}

\subsection{Mathias's criterion without a countability hypothesis}

Let $c=\seq{c_n:n<\omega}$ be strictly increasing and cofinal in
$\kappa$ in an outer universe containing $M$. Every finite initial
segment of $c$ belongs to $M$, since it is a finite sequence of
ordinals. Put
\begin{equation}\label{eq:Gc}
 G_c=\{(s,A)\in\PP_W^M:s\text{ is an initial segment of }c
          \text{ and all remaining entries of }c\text{ lie in }A\}.
\end{equation}

\begin{theorem}[Mathias genericity criterion]\label{thm:mathias}
The filter $G_c$ is $M$-generic if and only if
\begin{equation}\label{eq:Mathias}
  (\forall A\in W)\ (\exists k<\omega)\ (\forall n\ge k)\ c_n\in A.
\end{equation}
\end{theorem}

\begin{proof}
First, $G_c$ is a filter whenever \eqref{eq:Mathias} holds. It is
upward closed in the forcing order. For two conditions in it, their
stems are comparable initial segments of $c$. Take the longer stem
and intersect the two reservoirs above that stem. The intersection
belongs to $W$, contains the remaining tail of $c$, and gives a common
extension in $G_c$. Nonemptiness is immediate, for example from the
empty stem with reservoir $\kappa$.

Assume \eqref{eq:Mathias}, and fix a dense open $D\in M$. For every
finite stem $s$, apply Lemma~\ref{lem:strong-prikry} below the condition
with stem $s$ and full reservoir above it. By Choice in $M$, choose
simultaneously reservoirs $A_s\in W$ and integers $n_s$ such that all
$n_s$-point stem extensions from $A_s$, with their residual reservoir,
belong to $D$. Form $A=\Delta_{\fin} A_s\in W$ in $M$.

Choose $k$ so that $c_n\in A$ for all $n\ge k$, and set
$s=c\restr k$. For $n\ge k$, every entry of $s$ is less than $c_n$;
therefore $c_n\in A_s$ by the definition of the diagonal set. Let $t$
be the next $n_s$ entries of $c$. The strong Prikry conclusion gives
\[
 (s\concat t,A_s/t)\in D.
\]
All the still later entries of $c$ lie in this reservoir, so the
condition belongs to $G_c$. This proves $G_c\cap D\ne\varnothing$.
No enumeration of the dense subsets of the forcing was used.

Conversely, for any $A\in W$, conditions whose reservoirs are subsets
of $A$ form a dense set: intersect the current reservoir with $A$.
A generic filter meets that set. Beyond the corresponding stem, every
entry of its Prikry sequence is in $A$. This proves
\eqref{eq:Mathias}.
\end{proof}

\begin{corollary}[The Prikry subextension of a trace]\label{cor:factor}
Suppose $\NT(\lambda,b,q)$ holds, and write $N=N_{b,q}$,
$U=U_{b,q}$. Then every representative $a\in q$ is Prikry-generic over
$N$. The model $N[a]$ satisfies $\ZFC$, lies inside $V$, has
$V_\lambda^{N[a]}=V_\lambda$, and sees $\cf(\lambda)=\omega$.
For $a,d\in q$, $N[a]=N[d]$.
\end{corollary}

\begin{proof}
By definition of the trace, the increasing enumeration of $a$ is
eventually in every member of $U$. Apply Theorem~\ref{thm:mathias}
inside $N$, with this external sequence. Evaluation of $N$-names by the
filter \eqref{eq:Gc} takes place in $V$, giving a forcing extension
$N[a]\subseteq V$ satisfying $\ZFC$. It contains a cofinal increasing
$\omega$-sequence in $\lambda$.

The lower rank equality follows especially directly here:
$V_\lambda\subseteq N\subseteq N[a]\subseteq V$. Transitivity and
rank absoluteness give equality at height $\lambda$.

If $a=^*d$, their finite symmetric difference is a finite set of
ordinals, and so belongs to $N$. The operation of taking symmetric
difference shows $d\in N[a]$ and $a\in N[d]$. These mutual
containments imply equality of the two forcing extensions.
\end{proof}

This completes Theorem~\ref{thm:main}. It is useful to write $W_{b,q}$
for the common model $N_{b,q}[a]$, $a\in q$. This notation names a
Prikry subextension, not a claim that the full ambient universe is
$W_{b,q}$.

\subsection{Preservation facts needed for the converse construction}

We record the remaining classical facts with their proofs so that the
upper consistency bounds do not hide a preservation assumption.

\begin{proposition}[Prikry property and no bounded subsets]\label{prop:prikry-property}
Every sentence in the forcing language is decided by a direct
extension of any given condition. Consequently $\PP_W$ adds no new
subsets of any ordinal below $\kappa$.
\end{proposition}

\begin{proof}
Apply the strong Prikry lemma to the dense open set of conditions
deciding a fixed sentence $\varphi$. It gives $(s,A)$ and a length $n$
such that every canonical $n$-point stem extension decides
$\varphi$. Color these $n$-element stems by their two possible
decisions. Lemma~\ref{lem:ramsey} gives $B\subseteq A$ in $W$ on which
the decision is constant. Any condition below $(s,B)$ can be extended
to have at least $n$ new stem entries. It then extends one of the
canonical deciding conditions, and so makes the constant decision.
The conditions making that decision are dense below $(s,B)$; hence
$(s,B)$ itself decides $\varphi$. Its stem is still $s$.

Now let $\dot x$ be a name forced to be a subset of $\alpha<\kappa$.
Starting from any condition, use direct extensions to decide
$\xi\in\dot x$ successively for all $\xi<\alpha$. At limit stages,
intersect the previous reservoirs. There are fewer than $\kappa$
stages, so $\kappa$-completeness gives a reservoir in $W$, while the
stem never changes. The final direct extension decides all of
$\dot x$ as one ground-model subset of $\alpha$. Such conditions are
dense, so every generic extension has no new bounded subsets of
$\kappa$.
\end{proof}

\begin{proposition}[Cardinals and the lower rank segment]\label{prop:preservation}
If $G$ is $\PP_W$-generic over $M$, then in $M[G]$ all cardinals of
$M$ remain cardinals,
\[
 \cf(\kappa)=\omega,
 \qquad V_\kappa^{M[G]}=V_\kappa^M.
\]
\end{proposition}

\begin{proof}
The union of generic stems is increasing. Dense sets requiring a longer
stem make its order type $\omega$, and dense sets requiring an entry
above any specified $\beta<\kappa$ make it cofinal.

The no-bounded-subsets result preserves cardinals below $\kappa$:
a new collapse between two smaller ordinals could be coded by a new
subset of an ordinal below $\kappa$, using a ground-model pairing
bijection. The cardinal $\kappa$ itself remains a cardinal because it
was a limit of smaller cardinals, all still cardinals. A collapse of
$\kappa$ to an ordinal of size below it would also collapse one of
those intervening cardinals.

There are $\kappa$ possible finite stems, and two conditions with the
same stem are compatible by intersecting their reservoirs. Thus the
forcing is $\kappa^+$-cc. The usual antichain argument preserves all
cardinals at least $\kappa^+$: for a name for a function with a given
domain, each argument has at most $\kappa$ possible ordinal values.
A domain of ground-model cardinality $\nu$ can therefore cover at
most $\nu\cdot\kappa$ possible values, which cannot be a larger ground
cardinal.

Finally, $\kappa$ is inaccessible in $M$. Induct on $\alpha<\kappa$.
Assume the two models have the same $V_\alpha$. A ground-model
bijection of this set with an ordinal below $\kappa$ codes any of its
subsets by a bounded subset of $\kappa$. There are no new such codes,
so the models have the same $V_{\alpha+1}$. The limit step is union.
This proves equality at $\kappa$ as well.
\end{proof}

\section{Quotient-relative homogeneity and realization of traces}\label{sec:homogeneity}

We now reverse the construction. Start in a model $M\models\ZFC$ with
a normal measure $W$ on a measurable $\kappa$. Let $G$ be generic for
$\PP_W^M$, let $c$ be its increasing Prikry sequence, and work in
\[
 V^*=M[G],\qquad q=[\ran(c)].
\]
Choose in $M$ a bijection $b:\kappa\to V_\kappa^M$. This remains a
bijection onto $V_\kappa^{V^*}$ by Proposition~\ref{prop:preservation}.
The quotient $q$ in this section is computed in the extension $V^*$.

\subsection{Stem replacement fixes the quotient, not the sequence}

\begin{lemma}[Isomorphic cones with the same quotient]\label{lem:cones}
For any two conditions $(s,A)$ and $(t,B)$ there are direct extensions
$(s,C)$ and $(t,C)$ whose cones are isomorphic by
\[
 (s\concat u,D)\longmapsto(t\concat u,D).
\]
This isomorphism fixes ground-model parameters and preserves the
canonical name for the finite-difference class of the Prikry sequence.
\end{lemma}

\begin{proof}
Let $C=A\cap B$ restricted above every entry of $s$ and $t$. It belongs
to $W$. An extension below $(s,C)$ consists exactly of a finite new
stem $u$ from $C$ and a reservoir $D\subseteq C$ above the whole new
stem; replacing the old prefix $s$ by $t$ gives precisely an extension
below $(t,C)$. The inverse replaces $t$ by $s$, and the order is
preserved.

Paired generics for these cones give sequences $s\concat d$ and
$t\concat d$ with the same infinite tail $d$. Their ranges differ
by a finite set. They also generate the same extension of $M$: either
sequence is recovered from the other by replacing a ground-model
finite prefix. In this common extension the two finite-difference
classes are literally the same set. Thus, under the induced
translation of names, the names for $q$ are forced equal. Check-names
for ground-model sets, including $b$ and all ordinal parameters, are
fixed in the usual sense of cone isomorphisms.

This argument is a forcing-name argument and does not require actual
generics over the entire ambient universe. One can equivalently carry
it out in the Boolean completions of the two cones.
\end{proof}

\begin{lemma}[HOD containment despite the generic parameter]\label{lem:hod-ground}
In $V^*$,
\[
 K=\HOD_{\{b,q\}}^{V^*}\subseteq M.
\]
\end{lemma}

\begin{proof}
First consider any sentence with parameters $\check b$, the canonical
name $\dot q$, and check-names for ordinal parameters. It cannot be
forced true by one condition and false by another: shrink to the
isomorphic cones from Lemma~\ref{lem:cones}; their isomorphism preserves
the sentence and would carry a forcing of truth to a forcing of truth
below a condition forcing falsity. Since deciding conditions are dense,
the top condition decides every such sentence.

Suppose $x\subseteq\theta$ is ordinal definable from $b,q$ in $V^*$.
Fix a defining formula and its finite ordinal parameter tuple. For
each $\alpha<\theta$, the assertion that $\alpha$ belongs to the uniquely
defined set is a sentence of the preceding form. In $M$, use the
forcing relation to collect exactly the $\alpha$ for which the top
condition forces this assertion. Separation in $M$ produces a set
$x_0\subseteq\theta$, and the truth lemma shows $x=x_0$. Thus every
$\OD_{\{b,q\}}^{V^*}$ set of ordinals belongs to $M$.

To pass from sets of ordinals to hereditary definability, let $x\in K$.
In $V^*$, put the canonical $\OD_{\{b,q\}}$ order on $\tc(\{x\})$
and enumerate it in that order by an ordinal $\eta$. Code the membership
relation of this enumeration as a relation $E$ on $\eta$, and code $E$
as a set of ordinals using a fixed ordinal-pair coding. The code is
ordinal definable from $b,q$, so by the preceding paragraph $E\in M$.
It is well-founded and extensional. The well-foundedness witnessed in
$V^*$ implies the well-foundedness that $M$ needs to perform its
Mostowski collapse. The collapse is absolute and unique, and its value
at the ordinal corresponding to $x$ is $x$. Hence $x\in M$.
\end{proof}

The hereditary qualification is indispensable. The generic quotient
$q$ is itself definable from $q$, but it does not follow that $q\in M$.
The argument establishes containment for HOD, not for the entire class
of ordinal-definable objects from the quotient parameter.

\subsection{The old measure is recovered on the new inner power set}

\begin{theorem}[Prikry realization of a normal trace]\label{thm:realization}
In the above extension $V^*$, the triple $(\kappa,b,q)$ satisfies
$\NT$. For $K=\HOD_{\{b,q\}}^{V^*}$ its canonical trace is
\begin{equation}\label{eq:restriction}
 U_{b,q}=W\cap\Pow(\kappa)^K.
\end{equation}
If $W$ concentrates in $M$ on measurable cardinals below $\kappa$,
then $(\kappa,b,q)$ satisfies $\NTP$ in $V^*$.
\end{theorem}

\begin{proof}
The ambient cardinal and cofinality conditions follow from
Proposition~\ref{prop:preservation}. Lemma~\ref{lem:b-rank} gives
$V_\kappa^{V^*}\subseteq K$, whereas Lemma~\ref{lem:hod-ground} gives
$K\subseteq M$. In particular these models agree on $V_\kappa$.

For any $A\in\Pow(\kappa)^K$, we have $A\in M$. If $A\in W$, the
Prikry sequence is eventually in $A$, by genericity and the reservoir
density argument. If $A\notin W$, its complement belongs to $W$, so
the sequence is eventually outside $A$. Therefore
\[
 \ran(c)\subseteq^* A\quad\Longleftrightarrow\quad A\in W.
\]
This proves \eqref{eq:restriction}. Independently of that equality,
Lemma~\ref{lem:trace-set} puts the left side, as a set, inside $K$.
This is important: it is not being assumed that the original measure
$W$ itself belongs to $K$.

The right side of \eqref{eq:restriction} is a normal measure in $K$.
Indeed, a sequence in $K$ of length $\mu<\kappa$ of its members is a
sequence in $M$ of members of $W$. Its intersection belongs to $W$ by
completeness in $M$, and belongs to $K$ by the axioms of $K$, so it
belongs to the restricted measure. A regressive function in $K$ is
also in $M$. Normality of $W$ produces a constant fiber in $W$; its
ordinal value and the fiber belong to $K$, so normality holds for the
restriction. The ultrafilter and nonprincipal properties follow at
once from those of $W$ and equality of complements. Hence $\NT$ holds.

Finally, measurability of a cardinal $\alpha<\kappa$ is absolute among
$K,M,V^*$, because their $V_\kappa$ segments agree. Every measure
witness on $\alpha$, every subset of $\alpha$, and every sequence or
regressive function used to test it has rank below $\kappa$. Thus
\[
 \Meas_\kappa^K=\Meas_\kappa^M.
\]
If the latter set belongs to $W$, then it belongs to the trace by
\eqref{eq:restriction}. This proves $\NTP$.
\end{proof}

\begin{remark}[No definability assumption on the ground measure]
The construction works with an arbitrary normal measure in $M$.
We did not require $M=\HOD$, $W$ ordinal definable, or a definable
well-order of the whole universe. The ground-model bijection $b$ gives
all lower-rank objects to $K$, and the quotient $q$ defines the
restriction of $W$ to $K$.
\end{remark}

\section{The two equiconsistency theorems}\label{sec:calibration}

\subsection{Concentration and Mitchell height two}

For normal measures $W_0,W_1$ on $\kappa$, write
\[
 W_0\mathrel{\triangleleft}W_1
 \quad\Longleftrightarrow\quad W_0\in\Ult(V,W_1),
\]
where the ultrapower is identified with its transitive collapse. The
two-step form of ``$o(\kappa)\ge2$'' is the existence of
$W_0\triangleleft W_1$. This explicit formulation suffices throughout
and avoids any need for higher Mitchell-rank calculations.

\begin{lemma}[The concentration criterion]\label{lem:Mitchell}
For a normal measure $W$ on $\kappa$ in a universe of $\ZFC$, the
following are equivalent:
\begin{enumerate}
\item $\{\alpha<\kappa:\alpha\text{ is measurable}\}\in W$.
\item There is a normal measure $W_0$ on $\kappa$ with
$W_0\in\Ult(V,W)$.
\end{enumerate}
Thus existence of a measure as in (1) is equivalent to existence of a
cardinal of Mitchell order at least two.
\end{lemma}

\begin{proof}
Let $i:V\to M_W$ be the normal ultrapower. Normality gives
$[\id]_W=\kappa$: every ordinal predecessor represented by a regressive
function is equal modulo $W$ to a constant below $\kappa$. By \L o\'s's
theorem,
\begin{equation}\label{eq:LosMeas}
 \{\alpha<\kappa:\alpha\text{ is measurable}\}\in W
 \quad\Longleftrightarrow\quad
 M_W\models\text{``$\kappa$ is measurable.''}
\end{equation}

We verify that a normal measure witnessing the right side is also a
normal measure in the original universe. For every $A\subseteq\kappa$
in $V$,
\[
 A=i(A)\cap\kappa\in M_W.
\]
Consequently $M_W$ and $V$ have exactly the same subsets of $\kappa$.
They also have the same relations on $\kappa$ and the same sequences of
subsets of $\kappa$ of length at most $\kappa$: such a relation or
sequence is coded by a subset of $\kappa\times\kappa$, and a pairing
bijection in $M_W$ converts that code to a subset of $\kappa$.

If $M_W$ has a normal measure $W_0$ on $\kappa$, it therefore measures
all the ambient subsets, not a smaller power set. Every ambient
regressive function and every sequence of fewer than $\kappa$
measure-one sets used in testing normality or completeness is already
in $M_W$. The internal verification of these properties is thus an
ambient verification. This proves (1)$\Rightarrow$(2).

Conversely, if an ambient normal measure $W_0$ belongs to $M_W$, then
it is a normal measure there as well: the power sets agree, and the
internal tests are among the ambient tests. Hence the right side of
\eqref{eq:LosMeas} holds, proving (2)$\Rightarrow$(1).
\end{proof}

When this lemma is used in $N_{b,q}$, \emph{every occurrence of the
universe and ultrapower in it is interpreted in $N_{b,q}$}. It does not
turn the trace into a measure on the full ambient power set.

\subsection{Measurable strength for the basic trace}

\begin{proof}[Proof of \eqref{eq:conNT}]
For the lower bound, a model of $T_{\NT}$ has the explicitly definable
inner model $N_{b,q}$. By Lemma~\ref{lem:hod} this model satisfies
$\ZFC$, and the defining property of $\NT$ supplies its measurable
cardinal $\lambda$. Therefore consistency of $T_{\NT}$ implies
consistency of $\ZFC$ with a measurable cardinal.

For the upper bound, start with $\ZFC$ and a measurable $\kappa$,
choose a normal measure $W$, and force with $\PP_W$. Choose a
ground-model bijection $b:\kappa\to V_\kappa$ and let $q$ be the
finite-difference class of the generic sequence. Theorem~\ref{thm:realization}
proves that the extension satisfies $\NT(\kappa,b,q)$. Hence
consistency of the measurable-cardinal theory implies consistency of
$T_{\NT}$.
\end{proof}

\begin{corollary}[The stronger clause is not automatic for a trace]
The construction of Theorem~\ref{thm:realization} can be performed at
the first measurable cardinal of its ground model. At that particular
$\kappa$, $K=\HOD_{\{b,q\}}$ regards $\kappa$ as its first measurable
cardinal. Thus its trace is normal but does not concentrate on
measurables below $\kappa$.
\end{corollary}

\begin{proof}
The models $K$ and the ground agree below $\kappa$. No smaller
cardinal has a measure in the ground, and hence none does in $K$.
The canonical trace makes $\kappa$ measurable in $K$. The set of
smaller measurables is therefore empty and is not in the trace.
\end{proof}

This is a statement about the constructed triple. It is not a claim
that the entire resulting universe has no other possible $\NTP$
triple, which would require an additional global argument.

\subsection{Mitchell height two for the strengthened trace}

\begin{proof}[Proof of \eqref{eq:conNTplus}]
For the lower bound, a model of $T_{\NTP}$ has the inner model
$N_{b,q}\models\ZFC$ with a normal measure on $\lambda$ concentrating
on its measurable cardinals. Apply Lemma~\ref{lem:Mitchell} inside
that model. It sees $o(\lambda)\ge2$, establishing the required
consistency implication.

For the upper bound, assume there are normal measures
$W_0\triangleleft W_1$ on $\kappa$. By Lemma~\ref{lem:Mitchell},
$W_1$ concentrates on measurables below $\kappa$. Force with
$\PP_{W_1}$, choose the bijection $b$ in the ground, and form the
generic quotient $q$ in the extension. The last clause of
Theorem~\ref{thm:realization} gives $\NTP(\kappa,b,q)$ there.
\end{proof}

\subsection{Why these are ordinary consistency statements}

The model notation in the forcing proofs is the customary way to
present a relative consistency argument; it does not add the
hypothesis that a countable transitive model exists. The lower bounds
are uniform inner-model interpretations. For each axiom of $\ZFC$,
its relativization to the fixed-parameter HOD predicate is provable in
$\ZFC$, uniformly in the parameters. The measure assertions likewise
relativize by explicit first-order formulas.

For the upper bounds, $\ZFC$ with the stated measure hypothesis proves
that the appropriate set forcing forces every axiom of $\ZFC$ and the
respective first-order trace assertion. The forcing theorem, or its
Boolean-valued presentation, converts these proofs into the ordinary
syntactic relative consistency implication. Consequently
\eqref{eq:conNT} and \eqref{eq:conNTplus} concern $\Con(T)$, not the
stronger assertion $\Con(\text{there is a transitive model of }T)$.

\section{The HOD boundary and conditional incompatibilities}\label{sec:boundary}

\subsection{Eliminating the rank-bijection parameter}

Suppose the additional boundary condition $V_\lambda\subseteq\HOD^V$
holds. For $q\in Q_\lambda$ put
\[
 N_q^0=\HOD_{\{q\}}^V,\qquad
 U_q^0=\{A\in\Pow(\lambda)\cap N_q^0:
                      (\exists a\in q)\ a\subseteq^* A\}.
\]
The superscript $0$ distinguishes this model from the synthesis's
$\HOD_{V_\lambda\cup\{q\}}$ and from $N_{b,q}$.

\begin{theorem}[A quotient-only normal measure at the boundary]\label{thm:boundary}
Assume $\lambda$ is ultraexacting and $V_\lambda\subseteq\HOD^V$.
There is a critical-sequence class $q$ such that $N_q^0$ satisfies
$\ZFC$, has $V_\lambda^{N_q^0}=V_\lambda$, and regards $U_q^0$ as a
normal measure on $\lambda$ concentrating on measurables. Every
representative of $q$ is Prikry-generic over $N_q^0$ for $U_q^0$, and
all give the same Prikry subextension of $V$.
\end{theorem}

\begin{proof}
Repeat the height choice and critical-sequence construction, now using
only $p=q$ in the two counterexample predicates of
Lemmas~\ref{lem:decisions} and~\ref{lem:press}. The quotient is still
fixed by $j$. Those two proofs use no property of $b$ beyond its being
a fixed parameter, so they give eventual decisions and eventual
pressing down for all $\OD_{\{q\}}$ subsets and functions.

Lemma~\ref{lem:hod} gives $\ZFC$ for $N_q^0$. The boundary condition
supplies $V_\lambda\subseteq\HOD\subseteq N_q^0$, replacing the use
of Lemma~\ref{lem:b-rank}. The proof of Proposition~\ref{prop:normal}
therefore gives the normal measure, and the lower-rank measures on the
critical points give concentration as in
Proposition~\ref{prop:concentration}. Apply the Mathias criterion and
the finite-modification argument in Corollary~\ref{cor:factor} to
obtain all the remaining conclusions.
\end{proof}

\begin{corollary}[An enhanced I0-equiconsistent boundary package]\label{cor:boundary-con}
The following theory is equiconsistent with $\ZFC+\Izero$:
there is an ultraexacting $\lambda$ with $V_\lambda\subseteq\HOD$,
no exacting cardinal below $\lambda$, and a quotient-only normal
trace and Prikry subextension as in Theorem~\ref{thm:boundary}.
\end{corollary}

\begin{proof}
The published results \cite[Theorem~4.5, Proposition~3.9 and
Corollary~3.10]{abgl} supply, relative to $\Izero$, an ultraexacting
cardinal in a model with the displayed HOD boundary and no smaller
exacting cardinal. Theorem~\ref{thm:boundary} supplies the new normal
trace and Prikry clauses in that same model. Conversely, the stated
theory includes an ultraexacting cardinal, so its consistency implies
consistency of $\Izero$ by the published equiconsistency theorem.
\end{proof}

Unlike Theorem~\ref{thm:calibrations}, this last equiconsistency
\emph{inherits} its two large-cardinal bounds from the literature. Its
new content is the stronger structural conclusion in the boundary
model, not a new comparison between $\Izero$ and ultraexactingness.

\subsection{A precise conditional inconsistency}

\begin{corollary}[Low-measure HOD hypotheses are incompatible with
ultraexactingness]\label{cor:incompatibility}
There is no ultraexacting $\lambda$ satisfying the following additional
hypothesis:
\begin{quote}
For every bijection $b:\lambda\to V_\lambda$ and every $q\in Q_\lambda$,
$\HOD_{\{b,q\}}$ has no normal measure on $\lambda$ concentrating on
its measurable cardinals below $\lambda$.
\end{quote}
Equivalently, one may require all these models to have internal
$o(\lambda)\le1$ (with $o(\lambda)=0$ when there is no normal measure). At the HOD boundary, it suffices to quantify over
$\HOD_{\{q\}}$, without $b$.
\end{corollary}

\begin{proof}
Theorem~\ref{thm:main} contradicts the displayed hypothesis at its
specific pair $(b,q)$; Lemma~\ref{lem:Mitchell}, interpreted internally,
gives the equivalent Mitchell-order formulation. Under the boundary
condition, use Theorem~\ref{thm:boundary} instead.
\end{proof}

This is an explicit additional hypothesis about parameterized inner
models. It is not a theorem that this hypothesis holds in $\ZFC$, nor
an inconsistency of a standard large-cardinal axiom by itself. The
converse forcing constructions show why normal traces cannot be
rejected simply because their supporting sequences are countable in
the ambient universe.

\section{Internal completeness, omitted information, and proof audit}\label{sec:audit}

\subsection{The external countable intersection that does fail}

\begin{proposition}[Exactly where ambient completeness fails]\label{prop:external}
In the conclusion of Theorem~\ref{thm:main}, set
\[
 B_n=\lambda\setminus(\kappa_n+1)\qquad(n<\omega).
\]
Then every $B_n\in U_{b,q}$ but $\bigcap_{n<\omega}B_n=\varnothing$ in
$V$. The sequence $\seq{B_n:n<\omega}$ does not belong to $N_{b,q}$.
Moreover $q\notin N_{b,q}$ and $q\cap N_{b,q}=\varnothing$.
\end{proposition}

\begin{proof}
Every $B_n$ is a final segment, belongs to the inner model by its
ordinal definition, and eventually contains the critical sequence.
The intersection is empty because that sequence is cofinal in
$\lambda$. If the sequence of these sets were in $N_{b,q}$, internal
countable completeness would put the empty intersection in the
measure, contradicting propriety. Thus the sequence itself is missing
from the inner model.

Inside $N_{b,q}$, the cardinal $\lambda$ is measurable and therefore
regular. A representative $d\in q\cap N_{b,q}$ would have order type
$\omega$ and supremum $\lambda$ there as well: its increasing
enumeration and order type are absolute between the transitive models.
That would contradict regularity. Hence there is no such
representative. If $q$ itself belonged to the transitive model, so
would all of its members, which is impossible.
\end{proof}

There are two separate restrictions on the trace. Its domain is the
\emph{internal} power set $\Pow(\lambda)^{N_{b,q}}$, and its
completeness applies to sequences \emph{in} that model. Neither can
be replaced by its ambient version. In particular, the trace cannot
be inserted as an ambient complete ultrafilter into the HCD arguments
of the supplied completeness-barrier analysis.

The parameter $q$ can define the measure without belonging to the
hereditarily definable model. This is not a conflict with the HOD
construction: $U$ has only hereditarily definable descendants, whereas
$q$ has non-hereditarily-definable representatives among its members.

\subsection{A ground-model restriction, not a ground-model identification}

\begin{proposition}[A small chain-condition presentation is impossible]
For $N=N_{b,q}$ as in Theorem~\ref{thm:main}, $V$ is not a forcing
extension of $N$ by a forcing that is $\lambda$-cc in $N$.
\end{proposition}

\begin{proof}
Suppose such a forcing supplied a countable cofinal sequence in
$\lambda$. For each coordinate, choose in $N$ a maximal antichain
deciding its value. Each antichain has cardinality less than $\lambda$
in $N$. Their countable union has size less than $\lambda$ in $N$,
since $N$ regards $\lambda$ as regular uncountable. The set of all
possible values is therefore bounded below $\lambda$ in $N$, and the
same ordinal bound holds in $V$, a contradiction.
\end{proof}

This is the familiar regularity obstruction, now applied to the
normal-trace model. Our positive construction only proves
$N\subseteq N[a]\subseteq V$. It gives no forcing presentation of
$V$ over $N$ or over $N[a]$. Ordinary Prikry forcing is
$\lambda^+$-cc, not $\lambda$-cc, so the positive and negative
statements are compatible.

\subsection{Dependency and failure-mode audit}

\begin{longtable}{@{}p{.245\textwidth}p{.70\textwidth}@{}}
\toprule
\textbf{Step} & \textbf{What is used, and what is not assumed}\\
\midrule
\endhead
Correct height & A finite list of first-order predicates is reflected
uniformly in the later parameters. No full truth predicate for $V$ or
post hoc choice of a height for a previously fixed witness is used.\\[4pt]
Fixed objects & $e\in X$ puts the critical recursion and the union of
iterated bijections in $X$. Arbitrary elements of $X$ are not assumed
fixed.\\[4pt]
Canonical choice & A least $\OD_p$ counterexample is uniquely definable
from the fixed tuple. The ordinal parameters in its original
definition need not be fixed.\\[4pt]
Normality & The value of a fixed regressive function at the actual
critical point is below that point and hence fixed. This argument does
not apply to arbitrary orbit roots.\\[4pt]
Completeness & The least-failed-coordinate function is formed from a
sequence that belongs to the inner model. External sequences are
excluded, as Proposition~\ref{prop:external} demonstrates.\\[4pt]
Prikry genericity & The normal-measure diagonalization meets any one
inner-model dense set. No countability of the inner model is assumed.\\[4pt]
Homogeneity & Stem replacement preserves the entire finite-difference
class. It does not preserve a distinguished generic representative.\\[4pt]
Restriction of measure & The ground measure need not be in the HOD
model. Its restriction is there because it equals the canonically
definable trace.\\[4pt]
Mitchell order & The two-step calculation is performed in the named
universe. In the ultraexacting application it is an internal
$N_{b,q}$ calculation.\\[4pt]
Consistency & The two trace calibrations use inner-model
interpretations and set forcing. They do not use a transitive-model
existence hypothesis.\\[4pt]
Formal verification & No new Lean proof or Lean compilation is claimed.
The supplied Lean files were reference material, not a certificate for
the new normality or forcing arguments.\\
\bottomrule
\end{longtable}

\subsection{Questions that genuinely remain after this continuation}

The work does not settle whether ordinary exactingness rules out thin
definable complete sections, whether more than $\lambda$ rigid classes
must exist, or whether a cover-exacting cardinal can consistently lie
above a strongly compact cardinal. In particular, the external failure
of completeness prevents using the present measure as a shortcut
through the last question.

A focused new question is whether every rigid quotient class admits
\emph{some} canonical normal or extender trace, perhaps after a
change of the quotient presentation. The proof here does not establish
this: eventual pressing down uses the first point of the critical
sequence and its position below the critical point's image. Rigidity
alone records cardinalities of definable families and does not encode
that regressive-function calculation.

Another concrete question is whether the Prikry subextension
$W_{b,q}$ can be characterized without the chosen fixed bijection $b$
away from the HOD boundary. The boundary theorem removes $b$ under a
specific hypothesis, but no independence-of-$b$ theorem is proved in
the general case. These questions retain the structural information
exposed by the normal trace rather than trying to infer an ambient
measure on a singular cardinal.

\newpage
\section{Concluding perspective}

The critical sequence has two simultaneous interpretations. In $V$ it
is a countable cofinal sequence through a singular strong limit. In
its fixed-parameter HOD model, its finite-difference class determines
which internal subsets are eventually accepted, and the embedding
forces this acceptance rule to be a normal measure. Mathias's criterion
then identifies the same sequence as Prikry-generic over that model.

The converse construction explains the exact strength of the isolated
phenomenon. Finite stem changes are invisible to the quotient
parameter, so parameter-HOD stays inside a measurable ground model.
The generic quotient nevertheless reconstructs the old measure on the
inner power set. Adding concentration on measurables raises this
isolated consistency strength from a measurable cardinal to Mitchell
height two, and no further for the stated principle.

Thus the new conclusion is not an ambient contradiction. It is an
explicit normal-measure/Prikry structure beneath the ambient
singularity, together with two exact consistency calibrations and a
sharper conditional obstruction to uniformly low-measure
parameterized HOD models.

\phantomsection
\addcontentsline{toc}{section}{References}
\begin{thebibliography}{99}
\small
\setlength{\itemsep}{6pt}

\bibitem{source-synthesis}
\emph{Large cardinals at the inconsistency frontier: a synthesis}.
Second edition, 18 September 2026. Supplied in
\nolinkurl{Cardinals3.zip}, file
\nolinkurl{docs/Large_Cardinals_Synthesis.tex}.
Especially Section~12, ``Sections of the cofinal quotient.''

\bibitem{source-plan}
\emph{Large cardinals at the inconsistency frontier: a unified
assessment of ZFC-compatible hypotheses, their obstructions, and the
evidence on both sides}.
18 September 2026. Supplied in \nolinkurl{Cardinals3.zip}, file
\nolinkurl{docs/Large_Cardinals_Unified_Report.tex}.

\bibitem{abgl}
J.~P.~Aguilera, J.~Bagaria, G.~Goldberg, P.~L\"ucke.
\emph{Large cardinals beyond HOD}.
\arxiv{2509.10254v1}, 2025.
Definitions and height equivalences in Section~2;
Proposition~3.9, Corollary~3.10, and Theorem~4.5.

\bibitem{abl}
J.~P.~Aguilera, J.~Bagaria, P.~L\"ucke.
\emph{Large cardinals, structural reflection, and the HOD Conjecture}.
\arxiv{2411.11568v4}, 2025.

\bibitem{kunen}
K.~Kunen. Elementary embeddings and infinitary combinatorics.
\emph{Journal of Symbolic Logic} \textbf{36}(3) (1971), 407--413.
\doi{10.2307/2269948}.

\bibitem{mathias}
A.~R.~D.~Mathias. On sequences generic in the sense of Prikry.
\emph{Journal of the Australian Mathematical Society}
\textbf{15} (1973), 409--414.
\doi{10.1017/S1446788700028755}.

\bibitem{merimovich}
C.~Merimovich. \emph{Mathias like criterion for the extender based
Prikry forcing}. Author manuscript dated 12 October 2020.
The introduction recalls the classical criterion used here.
\url{https://www.carmimerimovich.net/publications/2021-Mathias_like_criterion_for_the_extender_based_Prikry_forcing.pdf}.

\end{thebibliography}
\end{document}
